\section{Metadata-Aware Architecture}
\label{sec:architecture}

\subsection{Architecture Overview}

The proposed system transforms a natural-language analytical request into an
executable Spark SQL query through a sequence of metadata-aware stages. The
architecture separates metadata selection, language-model generation,
validation, and physical query execution.

The complete processing flow is:

\begin{enumerate}

    \item receive a natural-language analytical question;

    \item retrieve metadata that is semantically relevant to the question;

    \item expand the initial metadata selection using semantic concepts,
    aliases, dataset evidence, and explicit cross-source relationships;

    \item render the selected metadata into a restricted SQL-safe context;

    \item ask a local language model to generate one Spark SQL query;

    \item parse and validate the generated SQL against the registered Spark
    schema and the expected output structure;

    \item if validation fails, optionally perform one repair attempt using the
    validation feedback and the same metadata context;

    \item execute the valid query using Spark SQL and return the computed
    result.

\end{enumerate}

The architecture is deliberately not implemented as a multi-agent system.
Deterministic Python orchestration is used for metadata processing,
validation, repair control, and Spark execution. The LLM is used only where
natural-language interpretation and SQL generation are required.



Figure~\ref{fig:system-architecture} summarizes the complete processing
pipeline and the separation between metadata grounding, generation,
validation, and execution.

\begin{figure}[htbp]
\centering

\begin{tikzpicture}[
    node distance=4mm,
    stage/.style={
        draw,
        rounded corners,
        align=center,
        text width=0.73\textwidth,
        minimum height=8mm,
        inner sep=4pt,
        font=\small,
        fill=gray!8
    },
    flow/.style={
        -{Latex[length=2.4mm]},
        thick
    }
]

\node[stage] (question)
    {Natural-language analytical question};

\node[stage, below=of question] (catalog)
    {Structured metadata catalog + local vector index};

\node[stage, below=of catalog] (retrieval)
    {Hybrid retrieval:
    dense Top-5 + lexical and alias evidence};

\node[stage, below=of retrieval] (expansion)
    {Semantic grounding + relation-aware expansion};

\node[stage, below=of expansion] (renderer)
    {SQL-safe metadata context};

\node[stage, below=of renderer] (llm)
    {Qwen2.5-Coder 3B:
    Text-to-Spark-SQL generation};

\node[stage, below=of llm] (validation)
    {SQLGlot + Spark validation
    $\rightarrow$ optional one-shot repair};

\node[stage, below=of validation] (spark)
    {Spark SQL execution
    $\rightarrow$ grounded analytical result};

\draw[flow] (question) -- (catalog);
\draw[flow] (catalog) -- (retrieval);
\draw[flow] (retrieval) -- (expansion);
\draw[flow] (expansion) -- (renderer);
\draw[flow] (renderer) -- (llm);
\draw[flow] (llm) -- (validation);
\draw[flow] (validation) -- (spark);

\end{tikzpicture}

\caption{
End-to-end metadata-aware Text-to-Spark-SQL architecture.
}
\label{fig:system-architecture}
\end{figure}

\subsection{Structured Metadata Catalog}

The metadata catalog is the central source of structural and semantic
knowledge used by the system. It is implemented as a relational SQLite
catalog and describes the curated Spark datasets rather than the original raw
files.

The core catalog contains four datasets and 66 physical columns. In addition
to the physical schema, it stores 21 semantic concepts, 29 column-to-concept
associations, 6 explicit relationships, 27 aliases, and 6 reusable semantic
rules.

The catalog separates several kinds of information:

\begin{itemize}

    \item \textbf{datasets}, including descriptions, physical paths,
    granularity, temporal information, and row counts;

    \item \textbf{columns}, representing the physical fields that can appear
    in executable Spark SQL;

    \item \textbf{semantic concepts}, providing a common vocabulary across
    heterogeneous physical schemas;

    \item \textbf{aliases}, connecting natural-language variants to datasets,
    concepts, and schema elements;

    \item \textbf{relationships}, describing executable join conditions
    between datasets;

    \item \textbf{semantic rules}, representing reusable SQL predicates or
    expressions associated with analytical concepts.

\end{itemize}

This separation allows semantic knowledge to be used during retrieval without
assuming that every metadata identifier is also a valid SQL identifier.


\subsection{Metadata Vector Index}

To support semantic retrieval, metadata entities are transformed into textual
documents and embedded using the local \texttt{embeddinggemma} model.

The base catalog produces 103 indexed metadata documents:

\begin{itemize}

    \item 66 column documents;
    \item 21 semantic-concept documents;
    \item 4 dataset documents;
    \item 6 relationship documents;
    \item 6 semantic-rule documents.

\end{itemize}

The resulting 768-dimensional embeddings are stored in a local Qdrant vector
index using cosine similarity.

The vector index is not used as an independent knowledge base. Its purpose is
to identify semantically promising metadata entities from the natural-language
question. Structural completeness is handled by the subsequent deterministic
expansion stage.


\subsection{Hybrid Metadata Retrieval}

The retrieval procedure combines dense semantic retrieval with deterministic
lexical and metadata-based evidence.

For each question, the system first retrieves the top five dense metadata
candidates from the vector index. Dense retrieval is complemented by lexical
matching against physical names and aliases. Explicit dataset mentions in the
question are treated as strong evidence and can override unrelated dataset
candidates returned only by embedding similarity.

Semantic concepts can also be expanded to their associated physical columns.
For example, a general concept such as pickup time can correspond to
\texttt{yellow\_taxi.tpep\_pickup\_datetime} or
\texttt{green\_taxi.lpep\_pickup\_datetime}, depending on the dataset
identified by the question.

The retrieval stage therefore does not simply select the five nearest vector
documents. Dense similarity is used as one signal inside a broader
metadata-grounding procedure.


\subsection{Relation-Aware Expansion}

Cross-source analytical queries often require metadata that is not directly
similar to the wording of the user question. A join key such as
\texttt{LocationID}, for example, may be essential for execution even if the
question mentions only a borough name.

To address this issue, the system performs deterministic relationship
expansion after the initial metadata retrieval. When selected datasets imply a
known catalog relationship, the relationship and its physical endpoint
columns can be added to the metadata selection.

Relationship selection also considers the query wording. This is necessary
when multiple relationships connect the same pair of datasets. Taxi records,
for example, can be related to Taxi Zones through either pickup or drop-off
location identifiers. Evidence in the question is therefore used to
distinguish pickup-oriented and drop-off-oriented relationships.

A similar mechanism is used for temporal relations between taxi pickups and
hourly weather observations. The catalog stores the executable relationship
between the hourly bucket derived from the taxi pickup timestamp and
\texttt{weather\_hourly.weather\_hour}.

This stage is one of the main differences between ordinary semantic retrieval
and the proposed relation-aware approach.


\subsection{Semantic Rules}

Some analytical concepts cannot be represented by selecting a physical column
alone. The catalog therefore contains semantic rules that associate common
analytical meanings with executable Spark SQL expressions or predicates.

For example, a rainy hour is represented by the physical predicate

\begin{verbatim}
weather_hourly.has_precipitation = TRUE
\end{verbatim}

while the taxi pickup hour is represented using an hourly transformation of
the corresponding pickup timestamp.

Rules are included in the retrieved context only when the selected metadata
and question provide evidence that they are relevant. This prevents the model
from receiving every semantic rule in the catalog for every request.


\subsection{SQL-Safe Context Rendering}

The metadata selected by retrieval cannot be passed directly to the language
model because the internal catalog contains identifiers that are meaningful
for metadata management but are not executable SQL.

The SQL-safe renderer therefore transforms the selected metadata into a
restricted prompt context. It exposes:

\begin{itemize}

    \item physical Spark table names;

    \item fully qualified physical column names;

    \item natural-language descriptions of their meaning;

    \item executable physical join conditions;

    \item executable SQL expressions and predicates.

\end{itemize}

Internal semantic-concept identifiers, relationship names, rule identifiers,
and aliases are not exposed as if they were physical SQL objects.

This distinction is important because the complete-catalog baseline contains
both physical and semantic metadata. Without SQL-safe rendering, an LLM can
incorrectly use an internal semantic identifier as a table or column name.

The renderer also provides dataset-level temporal information. Since the
registered taxi and weather datasets already contain only January 2024 data,
the context explicitly informs the model that an additional January predicate
is unnecessary unless the user requests a narrower date range.


\subsection{Text-to-Spark-SQL Generation}

The SQL generation component uses the local
\texttt{qwen2.5-coder:3b} model through Ollama \cite{hui2024qwen25coder}. Generation is deterministic,
with temperature set to zero and a context window of 4096 tokens.

The model receives the natural-language question together with either the full
catalog context or the relation-aware SQL-safe context, depending on the
experimental configuration.

Its task is restricted to generating a Spark SQL query. The model does not
directly compute the analytical answer. This ensures that final numerical
results are obtained from the physical datasets through Spark rather than from
the parametric knowledge of the language model.


\subsection{SQL Validation}

Generated SQL is validated before execution using a dedicated validation
stage.

The validator performs three main checks:

\begin{enumerate}

    \item parse the query using SQLGlot with the Spark SQL dialect;

    \item ensure that the generated statement is a single read-only query and
    that it references only allowed physical tables;

    \item ask Spark to analyse the query against the registered physical
    schemas and verify the expected output columns.

\end{enumerate}

This design catches syntactic errors, unknown tables, unknown physical
columns, and several forms of invalid schema grounding before actual query
execution.

Validation does not prove semantic correctness. A syntactically valid query
can still express the wrong analytical intention, which is why final
evaluation is based on query results rather than validation success alone.


\subsection{One-Shot SQL Repair}

When validation fails, the third system configuration allows one repair
attempt.

The repair model receives the invalid SQL, the validation feedback, the
original question, and the same retrieved metadata context used during the
initial generation. No gold SQL or expected answer is provided to the repair
stage.

Restricting repair to a single attempt keeps the experimental configuration
simple and prevents an uncontrolled iterative agent loop. If the repaired
query still fails validation, the system reports the failure instead of
continuing indefinitely.

The repair mechanism is intended to correct generation errors when the
required metadata is already available. It cannot recover information that
was omitted by the retrieval stage. This distinction becomes important in the
held-out error analysis.


\subsection{Spark SQL Execution and Grounded Results}

After successful validation, the query is executed using Apache Spark SQL over
the curated Parquet datasets.

Spark is therefore the authoritative computational layer of the system.
Aggregations, filters, joins, grouping operations, and result values are all
computed from the stored data.

The final architecture separates responsibilities clearly:

\begin{itemize}

    \item the metadata catalog represents structural and semantic knowledge;

    \item the retrieval layer selects the context required for a question;

    \item the LLM translates the grounded request into Spark SQL;

    \item validation checks whether the generated query is executable;

    \item optional repair handles a limited class of generation errors;

    \item Spark computes the final analytical result.

\end{itemize}

This separation also enables the experimental evaluation to analyse retrieval
quality, SQL correctness, LLM efficiency, catalog scalability, and Spark
data-volume scalability as distinct components of the overall system.
